\documentclass[a4paper,onecolumn,oneside,11pt,wide,floatssmall]{mwrep}
\usepackage[MeX]{polski}
\usepackage[utf8]{inputenc}
\usepackage{amssymb}
\usepackage[polish]{babel}

\let\emptyset\varnothing

\usepackage{bookman}
\usepackage[T1]{fontenc}

\usepackage{hyperref}
\usepackage{url}
\usepackage{breakurl}

\usepackage{wrapfig}
\usepackage{graphicx}
\usepackage{epstopdf}

\usepackage{units}

\usepackage{listings}
\lstset{language=C}
\renewcommand*\lstlistingname{Wydruk}
\renewcommand*\lstlistlistingname{Spis wydruków}


\usepackage[chapter]{algorithm}
\makeatletter
\renewcommand{\ALG@name}{Algorytm}
\renewcommand{\listalgorithmname}{Lista algorytmów}
\renewcommand{\thealgorithm}{\arabic{chapter}.\arabic{algorithm}.}
\makeatother

%% \usepackage[noend]{algorithmic}
\usepackage[noend]{algpseudocode}
%\algsetup{indent=2em}
\renewcommand{\algorithmicrequire}{\textbf{Wymaga:}}
\renewcommand{\algorithmicensure}{\textbf{Zapewnia:}}
\renewcommand{\algorithmiccomment}[1]{{\footnotesize \it \{ #1 \}}}
\renewcommand{\algorithmicif}{\textbf{Jeżeli}}
%% \renewcommand{\algorithmicthen}{\textbf{wówczas}}
\renewcommand{\algorithmicthen}{}
\renewcommand{\algorithmicelse}{\textbf{wpp.}}
%\renewcommand{\algorithmicelsif}{\algorithmicelse \textbf{jeżeli}}
\renewcommand{\algorithmicfor}{\textbf{Dla}}
\renewcommand{\algorithmicforall}{\textbf{Dla wszystkich}}
%% \renewcommand{\algorithmicdo}{\textbf{wykonuj}}
\renewcommand{\algorithmicdo}{\textbf{}}
\renewcommand{\algorithmicwhile}{\textbf{Dopóki}}
\renewcommand{\algorithmicloop}{\textbf{Wykonuj}}
\renewcommand{\algorithmicrepeat}{\textbf{Wykonuj}}
\renewcommand{\algorithmicreturn}{\textbf{Zwróć}}
\renewcommand{\algorithmicuntil}{\textbf{aż}}
\renewcommand{\algorithmicfunction}{\textbf{Funkcja}}
\renewcommand{\algorithmicprocedure}{\textbf{Procedura}}


\usepackage{geometry}
\textwidth\paperwidth
\advance\textwidth -55mm
\oddsidemargin-0.9in
\advance\oddsidemargin 33mm
\evensidemargin-0.9in
\advance\evensidemargin 33mm
\topmargin -1in
\advance\topmargin 25mm
\setlength\textheight{48\baselineskip}
\addtolength\textheight{\topskip}
\marginparwidth15mm

\clubpenalty=10000
\widowpenalty=10000
\brokenpenalty=10000
\sloppy

\tolerance4500
\pretolerance250
\hfuzz=1.5pt
\hbadness1450

\renewcommand{\chaptermark}[1]{\markboth{\scshape\small\bfseries \
#1}{\small\bfseries \ #1}}
\renewcommand{\sectionmark}[1]{\markboth{\scshape\small\bfseries\thesection.\
#1}{\small\bfseries\thesection.\ #1}}
\newcommand{\headrulewidth}{0.5pt}
\newcommand{\footrulewidth}{0.pt}
\pagestyle{uheadings}

\usepackage{color,graphicx}

\usepackage[sort,compress]{cite}

\newcommand{\TODO}[1]{{\setlength{\fboxsep}{0pt}\tiny (\colorbox{red}{TODO}: #1)}}
\newcommand{\ang}[1]{ang.\ {\it #1}}

\usepackage{titling}
\title{Alokacja ciągłych fizycznie obszarów pamięci w~systemie Linux}
\author{Michał Nazarewicz}
\date{\today}
\newcommand{\theengishtitle}{Physically contiguous memory alloction in Linux-based systems}


\begin{document}

\pagenumbering{roman}
\renewcommand{\baselinestretch}{1.0}
\raggedbottom

\begin{titlepage}
    % Strona tytułowa
    \vbox to\textheight{\hyphenpenalty=10000
    \begin{center}
        \begin{tabular}{p{107mm} p{9cm}}
            \begin{minipage}{9cm}
              \begin{center}
              Politechnika Warszawska \\
              Wydział Elektroniki i~Technik Informacyjnych \\
              Instytut Informatyki
              \end{center}
            \end{minipage}
            &
            \begin{minipage}{8cm}
            \begin{flushleft}
             \footnotesize
              Rok akademicki 2011/2012
            \vspace*{2.75\baselineskip}
            \end{flushleft}
            \end{minipage} \\
        \end{tabular}
        \vspace*{3.75\baselineskip}
        \par\vspace{\smallskipamount}
        \includegraphics[width=4cm]{build/pw.eps}\par
        \vspace*{2\baselineskip}{\LARGE Praca dyplomowa inżynierksa\par}
        \vspace{3\baselineskip}{\LARGE\strut \theauthor\par}
        \vspace*{2\baselineskip}{\huge\bfseries \thetitle\par}

        \vspace*{7\baselineskip}
        \hfill\mbox{}\par\vspace*{\baselineskip}\noindent
        \begin{tabular}[b]{@{}p{3cm}@{\ }l@{}}
            {\large\hfill } & {\large }
        \end{tabular}
        \hfill
        \begin{tabular}[b]{@{}l@{}}
        Opiekun pracy: \\[\smallskipamount]
        {\large dr inż.\ Wojciech Zabołotny}
        \end{tabular}\par
        \vspace*{4\baselineskip}
    \begin{tabular}{p{\textwidth}}
    \begin{flushleft}
        \begin{minipage}{7cm}
        Ocena \dotfill
        \par\vspace{1.6\baselineskip}
        \dotfill
        \par\noindent
        \centerline{\footnotesize Podpis Przewodniczącego} \par
        \centerline{\footnotesize Komisji Egzaminu Dyplomowego}\par
        \end{minipage}
    \end{flushleft}
    \end{tabular}
    \end{center}}

    % Życiorys
    \newpage\thispagestyle{empty}
    \begin{tabular}{p{5cm} p{12cm}}
    \begin{minipage}{5cm}
    \center
    \end{minipage}
    &
    \begin{minipage}{12cm}
    \begin{flushleft}
    \par\noindent\vspace{1\baselineskip}
    \begin{tabular}[h]{l l}
    {\it Kierunek:}       & Informatyka \\[12pt]
    {\it Specjalność:}    & Inżynieria Systemów \\
                          & Informatycznych \\[12pt]
    {\it Data urodzenia:} & 31 sierpnia 1986~r. \\[12pt]
    {\it Data rozpoczęcia studiów:} & 1 lutego 2006 r. \\
    \end{tabular}
    \par\noindent\vspace{1\baselineskip}
    \end{flushleft}
    \end{minipage}
    \end{tabular}
    \vspace*{1\baselineskip}
    \begin{center}
        {\large\bfseries Życiorys}\par\bigskip
    \end{center}

    \indent
    Nazywam się \theauthor \dots \TODO{do wypełnienia}
    \par
    \vspace{2\baselineskip}
    \hfill\parbox{15em}{{\small\dotfill}\\[-.3ex]
    \centerline{\footnotesize podpis studenta}}\par
    \vspace{3\baselineskip}
    \begin{center}
        {\large\bfseries Egzamin dyplomowy} \par\bigskip\bigskip
    \end{center}
    \par\noindent\vspace{1.5\baselineskip}
    Złożył egzamin dyplomowy w dn. \dotfill
    \par\noindent\vspace{1.5\baselineskip}
    Z wynikiem \dotfill
    \par\noindent\vspace{1.5\baselineskip}
    Ogólny wynik studiów \dotfill
    \par\noindent\vspace{1.5\baselineskip}
    Dodatkowe wnioski i uwagi Komisji \dotfill
    \par\noindent\vspace{1.5\baselineskip}
    \dotfill

    % Streszczenie
    \newpage\thispagestyle{empty}
    \vspace*{2\baselineskip}
    \begin{center}
        {\large\bfseries Streszczenie}\par\bigskip
    \end{center}

    {\itshape Praca ta prezentuje sposób alokacji dużych obszarów
      ciągłej fizycznie pamięci w~systemach opartych na jądrze Linux.
      Zastosowanie w~procesorach jednostek zarządzania pamięcią
      pozwala na uniknięcie konieczności takich alokacji, jednak
      w~systemach wbudowanych poszczególne komponenty są często
      pozbawiony jednostki translacji adresów pamięci, przez co muszą
      operować bezpośrednio na adresach fizycznych.  Aby temu
      zaradzić, stworzyłem alokator pamięci ciągłej (ang. {\it
        Contiguous Memory Allocator}, CMA), który pozwala przydzielać
      duże ciągłe fizycznie obszary pamięci, bez konieczności
      rezerwowania na wyłączność dużych obszarów przy starcie systemu.}
    \vspace*{1\baselineskip}

    \noindent{\bf Słowa kluczowe}: {\itshape alokacja pamięci, Linux,
      systemy wbudowane.}
    \par
    \vspace{4\baselineskip}
    \begin{center}
        {\large\bfseries Abstract}\par\bigskip
    \end{center}
    \noindent{\bf Title}: {\itshape \theengishtitle}\par
    \vspace*{1\baselineskip}

    {\itshape This thesis describes a~way to allocate large physically
      contiguous memory regions in Linux-based systems.  Memory
      management units allows to avoid the need of such allocations,
      but in embedded systems, individual components often lack an
      address translation unit and thus need to operate directly on
      physical addresses.  To address this issue, I~have created
      a~Contiguous Memory Allocator (or CMA) which allows to allocate
      big physically contiguous memory regions without the need to
      reserve memory for exclusive use during system start.}
    \vspace*{1\baselineskip}

    \noindent{\bf Key words}: {\itshape memory allocation, Linux,
      embedded systems.}

\end{titlepage}



\tableofcontents

% \chapter*{Spis tablic, rysunków i~wydruków}
% \listoftables
% \listoffigures
% \lstlistoflistings


\newpage
\pagenumbering{arabic}
\setcounter{page}{1}

\chapter{Wstęp}

Nowoczesne procesory przeznaczone do serwerów i~komputów osobistych
posiadają jednostkę zarządzania pamięcią (ang.\ {\it memory management
  unit}, MMU), która tłumaczy adresy logiczne, widziane przez
procesor, na adresy fizyczne.  Dzięki temu, obszary pamięci, które
z~punktu widzenia procesora (a~więc programów na nim uruchomionych)
mają ciągłe adresy, mogą w~rzeczywistości być podzielone na wiele
stron rozrzuconych po pamięci fizycznej.  W~ten sposób, nawet jeżeli
program alokuje wielomegabajtowy\footnote{Aby uniknąć wieloznaczności,
  stosuję przedrostki „kilo-” i~„mega-” w~znaczeniu odpowiednio tysiąc
  i~milion (zgodnie z~układem SI), a~„kibi-” i~„mebi-” w~znaczeniu
  odpowiednie $2^{10} = 1024$ i~$2^{20}$ (zgodnie
  z~tzw.\ przedrostkami ICE).} obszary, system może zrealizować żądanie
alokując wiele czterokilobajtowych stron i~nie przejmować się
fragmentacją pamięci.

Niemniej, MMU nie jest zazwyczaj dostępny dla pozostałych układów
zintegrowanych znajdujących się w~urządzeniu, takich jak np.\ karta
dźwiękowa, czy kontroler sieciowy.  Jednak także i~dla tych przypadków
istnieją rozwiązania w~postaci mechanizmu bezpośredniego dostępu do
pamięci (ang.\ {\it Dircet Memory Access}, DMA) z~obsługą mechanizmu
vectored IO \TODO{przetłumaczyć}, który pozwala na „złożenie” buforu
z~wielu mniejszych buforów (stąd też inna nazwa tego mechanizmu:
ang.\ {\it scatter-gather}).

Niestety, zarówno MMU jak i~DMA obsługujące vectored IO
\TODO{przetłumaczyć} nie są pozbawione kosztów -- zarówno związanych
z~koniecznością fizycznego wyprodukowaniom odpowiedniego układu, ale
także związanych z~szybkością dostępu do pamięci.  Z~tego powodu,
w~systemach wbudowanych, takich jak np.\ telefony komórkowe, takie
mechanizmy są często niedostępne dla zewnętrznych (z~punktu widzenia
procesora) układów.  Jednocześnie, właśnie takie urządzenia posiadają
dużo takich układów, jak chociażby aparat fotograficzny, czy układ do
szybkiego dekodowania strumieni multimedialnych.

Powoduje to, że tego typu układy muszą operować bezpośrednio na
adresach fizycznych i~w~konsekwencji, wszelkie stosowane przez nie
bufory muszą być ciągłe w~pamięci fizycznej.  Niestety, Linux nie jest
dobrze przystosowany do alokowania dużych obszarów ciągłych
fizycznie\TODO{zbyt dużo „fizycznie”}\footnote{Dla przykładu,
  pięciomegapikselowa kamera potrzebuje buforu o~rozmiarze 15
  megabajtów, a~pojedyncza ramka {\it full HD} (tj.\ $1920 \times
  1080$) zajmuje ponad sześć megabajtów (i~dodatkowo dekoder musi
  często przechowywać kilka ramek).  Linux nie jest nawet w~stanie
  (bez modyfikowania źródła) zarządzać obszarami większymi niż cztery
  megabajty (1024 strony).  Na domiar złego fragmentacja zewnętrzna
  utrudnia alokacje, a~fragmentacja wewnętrzna powoduje duże straty
  pamięci.}.

Najprostszym, i~stosunkowo często stosowanym, rozwiązaniem tego
problemu jest zarezerwowanie przy starcie systemu pewnego obszaru
pamięci na potrzeby danego sterownika.  Niestety, o~ile taki mechanizm
może być wystarczający, jeżeli podzespoły wymagają stosunkowo małych
buforów, przestaje się on skalować przy współczesnych systemach, gdyż
wymaga rezerwacji wielu megabajtów pamięci, która przez większość
czasu nie jest do niczego wykorzystywana.

\section{Physical Memory Manager}\label{sec:pmm}

Bardziej skomplikowanym rozwiązaniem jest mechanizm, który rezerwuje
pewną przestrzeń pamięci, ale zamiast na stałe przypisywać regiony
pamięci do urządzeń, pozwala sterownikom alokować bufory, wtedy, gdy
są one potrzebne.

W~trakcie moich prac napisałem {\it Physical Memory Manager} (lub PMM)
\cite{patch:pmm}, który implementuje dokładnie te założenia. W~tym
podstawowym założeniu, PMM nie przedstawia sobą nic nowego.  Bowiem
w~1996 roku Matt Welsh napisał pierwszą wersję patcha
\emph{bigphysarea} dla jądra 1.3.71, który był z~różnym zaangażowaniem
utrzymywany i~przystosowywany do nowszych wersji Linuksa
\cite{patch:bigphys}.

PMM umożliwiał alokowanie dużych obszarów ciągłej pamięci fizycznej
nie tylko sterownikom działającym w~przestrzeni jądra, ale także
programom działających pod kontrolą systemu.  W~ten sposób, aplikacja
mogła zaalokować pamięć, tak aby (dla przykładu) dekoder JPEG mógł
zdekodować plik graficzny i~umieścić go w~pamięci bezpośrednio
dostępnej dla aplikacji.  W~ten sposób, kod odpowiedzialny za
mapowanie zaalokowanej pamięci był umieszczony wewnątrz PMM in
poszczególne sterowniki nie musiały się tym problemem zajmować.

Co więcej, PMM został zintegrowany z~mechanizmem współdzielenia
pamięci Systemu V~(tj.\ funkcjami \lstinline|shmget()|,
\lstinline|shmat()|, \lstinline|shmdt()| itp).  Mechanizm ten jest
wykorzystywany między innymi przez system X~Window do umożliwienia
współdzielenia pixmap pomiędzy klientem i~serwerem (działającymi na
tej samej maszynie) bez konieczności przysłania danych przez gniazdo
sieciowe.

Pozwalało to na dekodowanie obrazów i~strumieni video bezpośrednio do
obszarów pamięci z~których X11 odczytywał dane.  Dzięki temu, w~całym
procesie dane nie były niepotrzebnie kopiowane z~jednego miejsca
w~pamięci do drugiego, co minimalizowało zużycie procesora i~szyny
pamięci i~w~rezultacie przyśpieszało działanie systemu.

Jednakże, pamięć zarezerwowana przez PMM i~tak przez większość czasu
jest zupełnie nieużywana, a~więc marnowana.  Z~tego powodu, PMM
w~formie zaprezentowanej na początku nie został przyjęty przez
społeczność programistów Linuksa i~musiałem rozwijać go dalej.

\section{Contiguous Memory Allocator}\label{sec:cma}

Aby rozwiązać i~ten problem, zacząłem rozwijać {\it Contiguous Memory
  Allocator} (lub CMA).  W~swoich początkowych wersjach, i~ten
mechanizm działał na założeniach identycznych do PMM -- rezerwował
pamięć przy starcie systemu, którą potem zarządzał pozwalając
sterownikom i~programom na alokowanie obszarów ciągłych fizycznie
\cite{patch:cma-4}.

Pierwsze wersje CMA skupiały się głównie na rozwiązywaniu problemu
przypisywania różnych zarezerwowanych obszarów do różnych urządzeń,
a~także umożliwianiu sterownikom alokowanie różnych buforów w~różnych
obszarach pamięci.  Było to potrzebne, gdyż dekoder wideo (ang.\ {\it
  Multi-Format Codec}, lub MFC) stosowany na platformie S5PV110
wymagał, aby różne dane były dostępne w~różnych bankach pamięci.
Pozwalało to na zwiększenie szybkości dostępu do tych danych dzięki
zastosowaniu odczytu z~dwóch banków pamięci jednocześnie.

Z~czasem, coraz bardziej integrowałem CMA z~kodem zarządzania pamięci
w~Linuksie w~wyniku czego, pamięć rezerwowana przy starcie systemu,
stała się dostępna dla reszty systemu, jeżeli żaden sterownik jej nie
używał \cite{patch:cma-24}.  Takie rozwiązanie zostało ostatecznie
zaakceptowane przez społeczność deweloperów Linuksa i~jest dostępne
w~źródłach Linuksa począwszy od wersji 3.5.  W~tej pracy opisuję CMA
w~formie w~jakiej znalazł się on w~Linuksie 3.5\footnote{Należy
  zauważyć, że Linuks jest szybko rozwijającym się projektem wolnego
  oprogramowania i~ponieważ CMA używana jest przez coraz więcej osób,
  jest ona ciągle rozwijana i~im dalej w~przyszłość, tym bardziej opis
  w~niniejszej pracy będzie się różnił od stanu faktycznego.}.



\appendix

\nocite{*}
\bibliographystyle{plplain}
\begin{thebibliography}{9}

\bibitem{bib:ldd3}
  J. Corbet, A. Rubini i G. Kroah-Hartman,
  \emph{Linux Device Drivers, Third Edition}.
  2005.
  \url{http://lwn.net/Kernel/LDD3/}

\bibitem{bib:taocp-fa}
  Donald Knuth,
  \emph{The Art of Computer Programming. Volume 1: Fundamental
    Algorithms},
  Third Edition,
  Massachusetts, 1997.
  ISBN 0-201-89683-4

\bibitem{bib:compaction}
  J. Corbet,
  \emph{Memory compaction},
  6 stycznia 2010.
  \url{http://lwn.net/Articles/368869/}

\bibitem{bib:cma-and-arm}
  J. Corbet,
  \emph{CMA and ARM},
  5 czerwca 2011.
  \url{http://lwn.net/Articles/450286/}

\bibitem{patch:cma-24}
  M. Nazarewicz i~M. Szyprowski,
  \emph{Contiguous Memory Allocator, v.\ 24},
  3 kwietnia 2012.
  \url{http://thread.gmane.org/gmane.linux.kernel.mm/76241}

\bibitem{patch:pmm}
  M. Nazarwicz,
  \emph{Physical Memory Management}.
  13 maja 2009.
  \url{http://lkml.org/lkml/2009/5/13/100}

\bibitem{patch:cma-4}
  M. Nazarewicz,
  \emph{The Contiguous Memory Allocator framework, v.\ 4},
  20 sierpnia 2010.
  \url{http://thread.gmane.org/gmane.linux.kernel.mm/52010}

\bibitem{patch:bigphys}
  D. Tochansky,
  \emph{bigphysarea patch for 3.2.x},
  27 marca 2012.
  \url{http://article.gmane.org/gmane.linux.kernel/1273100}

\bibitem{patch:cma-test}
  Barry Song,
  \emph{A simple kernel module as a helper to test CMA, v. 4},
  7 marca 2012.
  \url{http://thread.gmane.org/gmane.linux.kernel/1263136}

%% \bibitem{bib:state-of-allocators}
%%   D. Rientjes,
%%   \emph{Status of the Linux Slab Allocators}.
%%   Southern California Linux Expo, February 26, 2011
%%   \url{http://www.socallinuxexpo.org/scale9x/presentations/status-linux-slab-allocators.html}

\end{thebibliography}


\end{document}
